\section{Robustness}\label{sec:10_robustness}

The mechanism rests on three facts: the entry floor, the trading bounds of \Cref{sec:05_investor-bounds}, and
recruitment after a purchase. Each extension below changes one primitive and asks which facts survive.
\paper{Sec.~6; Appendix~A.6} The bounds \(\rho\kappa\DT\le K_H,K_L\le\DT\) still concern only the gross returns of
correctly signed nonzero orders.

\subsection{Preparation cost laws: keep the floor first}\label{sec:10_robustness-costs}
\paper{Sec.~6.1; Prop.~A.4}

\begin{intuition}[The atoms do not carry the result]
The zero-cost type prepares after every price, so entry never falls below \(\rho\). The high-cost type prepares only
after good news. Any cost law with a low and a high component can do both jobs, if strict gaps separate the costs
from the gross profits.
\end{intuition}

Let the cost be independent of the class, the values and the slot, with law \(H_C=\rho H^0+(1-\rho)H^1\),
\(\operatorname{supp}H^0\subset[0,L]\) and \(\operatorname{supp}H^1\subset[c_-,c_+]\). Either component may have atoms.

\begin{resultbox}[Preparation cost laws \paper{Prop.~A.4}]
Let \(F_1\) strengthen \(F_0\), and suppose
\[
0<L<B_1(m),\qquad B_0(\pi)<c_-<c_+<B_1(M),\qquad \Delta_0<k<\rho\kappa\Delta_1.
\]
At \(F_0\), orders are uniquely zero, and entry is \(\rho\). At \(F_1\), orders are uniquely full, only a purchase
recruits the high component, and entry is \(\rho+(1-\rho)(b+\pi\lambda)\). Ownership rises under the exact condition
of Proposition~2.
\end{resultbox}

The proof of Proposition~2 goes through with three changes.
\begin{enumerate}[leftmargin=2em]
\item \emph{The floor.} In any candidate at \(F_0\) or \(F_1\), every price belief lies in \([m,M]\), also at a pooled price
(\Cref{sec:04_slot_beliefs-ceiling}; a pool with the zero order uses \(\beta\le1/3\)). Since \(g_H\ge g_L\) and
\(B_0\ge B_1\), \(B_i(\bar\mu)\ge B_i(m)\ge B_1(m)>L\). The whole low component prepares after every price, so
\(e\ge\rho\).
\item \emph{The orders.} The trading bounds need only this floor. At \(F_0\) a correct order earns at most
\(\Delta_0-k<0\); at \(F_1\) at least \(\rho\kappa\Delta_1-k>0\). So orders are zero at \(F_0\), where
\(B_0(\pi)<c_-\) keeps the high component out, and full at \(F_1\).
\item \emph{Recruitment.} \(B_1(M)>c_+>c_->B_0(\pi)\ge B_1(\pi)\) and \(M>\pi\) give \(g_{H,1}>g_{L,1}\). So
\(B_1(m)<B_1(\pi)<c_-\) keeps the high component out after a sale or zero, and \(B_1(M)>c_+\) admits all of it after a
purchase.
\end{enumerate}
Every cost comparison is strict, so no tie rule matters. Entry is at least \(\rho\) and nondecreasing in the
posterior, so the price algebra of \Cref{sec:04_slot_beliefs-sufficiency} applies unchanged, and entry and ownership
follow as in \Cref{sec:06_main_result-entry}.

\paragraph{The third economy needs another gap.}
At \(F_2\), the conditions \(L<B_2(m)\), \(B_2(M)<c_-\) and \(k<\rho\kappa\Delta_2\) keep full orders and return entry
to \(\rho\). Without \(B_2(M)<c_-\), part of the high component still prepares after a purchase. With full orders and
\(B_r(\pi)<c_-\), entry is \(\Ent(r)=\rho+(1-\rho)(b+\pi\lambda)\,H^1\bigl(B_r(M)\bigr)\). Only here does atomlessness
matter: if \(H^1\) has no atom, entry is continuous and nonincreasing in \(r\), and the drop at the entry ceiling
becomes a ramp.

\paragraph{Benchmark check.}
Take \(L=\reg{atom_L}\) and \(H^1\) uniform on \([6-\varepsilon,6+\varepsilon]\). At \(\varepsilon=\reg{atom_eps}\),
the support margins are \(B_1(m)-L=\reg{margin_p4_atomless_floor_B}\), \(c_--B_0(\pi)=\reg{margin_p4_atomless_weak_B}\)
and \(B_1(M)-c_+=\reg{margin_p4_atomless_strong_B}\), and entry at \(r_1\) is the two-point value
\(\reg{p4_atomless_r1_entry}\). At \(r_2\), \(B_2(M)=\reg{p2_base_r2_low_profit_plus}\) lies inside the band, and
\(H^1(B_2(M))=(B_2(M)-6+\varepsilon)/(2\varepsilon)\) gives \(\Ent_2=\reg{p4_atomless_r2_entry}\) and
\(\Own_{H,2}=\reg{p4_atomless_r2_own}\). Entry returns to \(\rho\) at \(r_2\) for every half-width below
\(6-B_2(M)=\reg{p4_atomless_max_width_floor_r2}\), such as \(\reg{atom_eps_small}\).

\begin{caution}[Keep the three parts separate]
The half-width \(\reg{atom_eps}\) keeps the rise from \(r_0\) to \(r_1\), but not entry \(\rho\) at \(r_2\).
Atomlessness smooths the drop at the ceiling; it does not remove it.
\end{caution}

\subsection{General value classes and any prior}\label{sec:10_robustness-classes}
\paper{Sec.~6.2; Prop.~A.5; Table~A.4}

\begin{intuition}[Put both strengths just above the point where the classes split]
Near the lowest value at which the classes differ, the target-payoff spread is close to zero, so informed trade
against the weaker incumbent cannot cover \(k\). The gross-profit gap there is close to the whole class gap, so a
high cost that only good news justifies fits inside it.
\end{intuition}

Let the classes satisfy the class ranking with \(\E_H[\theta]>\E_L[\theta]\). With \(D=G_L-G_H\ge0\), define
\(\Gamma(x)=\int_p^xD(y)\,dy\), \(\bar\Gamma=\E_H[\theta]-\E_L[\theta]>0\) and \(v_*=\inf\{x\ge p:\Gamma(x)>0\}\).
\(\Gamma\) rises from \(0\) to \(\bar\Gamma\), and \(v_*\) is the first value above the reserve where the classes differ.
For point masses, \(v_*=\ell\).

\emph{Step 1: the spread vanishes at \(v_*\).} For \(R\sim U[0,r]\), the spread formula \paper{(8)} has \(D(p^-)=0\) and
\(1-F(y)=1-y/r\). Integrate by parts with \(\Gamma'=D\) and \(\Gamma(p)=0\); the boundary term is zero at both ends:
\[
\DT(r)=\int_p^r\Bigl(1-\frac yr\Bigr)D(y)\,dy
=\Bigl[\Bigl(1-\frac yr\Bigr)\Gamma(y)\Bigr]_p^r+\frac1r\int_p^r\Gamma(y)\,dy
=\frac1r\int_p^r\Gamma(y)\,dy .
\]
Since \(\Gamma\) rises, \(\int_p^r\Gamma\le(r-p)\Gamma(r)\), so
\(\DT'(r)=[r\Gamma(r)-\int_p^r\Gamma]/r^2\ge p\Gamma(r)/r^2>0\) for \(r>v_*\). Since \(\Gamma=0\) on \([p,v_*]\),
\(\DT(r)\to0\) as \(r\downarrow v_*\).

\emph{Step 2: the profit gap stays large.} The surplus split of \Cref{sec:03_auction-split} gives
\(\DT+g_H-g_L=\bar\Gamma\) at every strength. With \(\DT(r)\le(r-p)\bar\Gamma/r\) from Step 1,
\(g_H(r)-g_L(r)=\bar\Gamma-\DT(r)\ge p\bar\Gamma/r>0\). At \(r=v_*\) the gap is the whole \(\bar\Gamma\), so
\(B_{v_*}(M)-B_{v_*}(\pi)=(M-\pi)\bar\Gamma>0\). Also \(W_H(r)\ge p/r>0\), since every value is at least \(p\), and
\(g_j\) and \(W_H\) are continuous in \(r\).

\begin{resultbox}[Construct the strengths and the costs in this order \paper{Prop.~A.5}]
Let \(e_H=\rho+(1-\rho)(b+\lambda)>\rho\).
\begin{enumerate}[leftmargin=2em]
\item Choose \(r_1>v_*\) close to \(v_*\), with \(B_{r_1}(M)>B_{v_*}(\pi)\) and \(e_HW_H(r_1)>\rho W_H(v_*)\). Both
hold at \(r_1=v_*\), so continuity permits the choice.
\item Choose \(r_0\in(v_*,r_1)\) close to \(v_*\), with \(\DT(r_0)<\rho\kappa\DT(r_1)\). Gross profit and \(W_H\) fall
with strength, so \(B_{r_0}(\pi)<B_{r_1}(M)\) and \(e_HW_H(r_1)>\rho W_H(r_0)\).
\item Choose \(c_H\in\bigl(B_{r_0}(\pi),B_{r_1}(M)\bigr)\) and \(k\in\bigl(\DT(r_0),\rho\kappa\DT(r_1)\bigr)\).
\item Since \(g_j(r)\to0\) as \(r\to\infty\), choose \(r_2>r_1\) with \(B_{r_2}(M)<c_H\).
\end{enumerate}
Every condition is a strict inequality between continuous functions. So, for every prior, reserve, \(\rho\) and slot,
a nonempty open set of \((r_0,r_1,r_2,c_H,k)\) satisfies (H1)--(H3) and the ownership condition. For point masses this
holds for every \(h>\ell\).
\end{resultbox}

\paragraph{Six economies.}
Each row keeps the benchmark slot, \(p=\reg{prim_p}\) and \(\rho=\reg{prim_rho}\), with its own costs and strengths. A
single value is a point mass; a bracket is a uniform class. In every row, entry at \(r_0\) is \(\rho\), and so is
entry at \(r_2\) when the row has a third strength.
\begin{center}\small
\setlength{\tabcolsep}{3pt}
\begin{tabular}{@{}lccccccc@{}}\toprule
Economy & Low class & High class & \(\pi\) & \((c_H,k)\) & \((r_0,r_1,r_2)\) & \(\Ent_1\) & \(\Own_{H,1}\)\\\midrule
Moderate gap & \(1\) & \(2\) & \(1/2\) & \((89/100,1/500)\) & \((21/20,3/2,\text{--})\)
 & \reg{p2_moderate_r1_entry} & \reg{p2_moderate_r1_ownership}\\
Narrow bands & \([19/20,21/20]\) & \([199/20,201/20]\) & \(1/2\) & \((6,1/50)\) & \((6/5,3,18/5)\)
 & \reg{p2_bands_r1_entry} & \reg{p2_bands_r1_ownership}\\
Low prior & \(1\) & \(10\) & \(3/10\) & \((9/2,1/100)\) & \((11/10,3,18/5)\)
 & \reg{p2_ex_a_r1_entry} & \reg{p2_ex_a_r1_ownership}\\
High prior & \(1\) & \(10\) & \(7/10\) & \((36/5,1/100)\) & \((11/10,3,18/5)\)
 & \reg{p2_ex_b_r1_entry} & \reg{p2_ex_b_r1_ownership}\\
Wide classes & \([3/5,7/5]\) & \([1,19]\) & \(7/10\) & \((67/10,1/50)\) & \((1,9/2,11/2)\)
 & \reg{p2_ex_c_r1_entry} & \reg{p2_ex_c_r1_ownership}\\
Same means & \(1\) & \(10\) & \(7/10\) & \((67/10,1/50)\) & \((1,9/2,11/2)\)
 & \reg{p2_ex_d_r1_entry} & \reg{p2_ex_d_r1_ownership}\\\bottomrule
\end{tabular}
\end{center}
The first five rows show the rise; the moderate-gap row has no third strength, so it checks (H1) and (H2) only. The
last row puts point masses at the means of the wide classes: \(B_1(M)-c_H=\reg{margin_p2_ex_d_strong_top_B}<0\), so a
purchase does not recruit. The prior rows need their own costs: at the benchmark costs and strengths, the prior
\(3/10\) gives entry \(\reg{p2_fixed_prior_three_tenths_r1_entry}\) at every strength, and the prior \(7/10\) gives
entry \(\reg{p2_fixed_prior_seven_tenths_r0_entry}\) at \(r_0\), because the prior itself recruits.

\emph{How to compute a row.} For a point value \(v\), the payoff terms of \Cref{sec:03_auction-terms} apply with \(v\)
in place of \(\ell\) when \(p\le v\le r\) (\(W_v=v/r\)), and with \(v\) in place of \(h\) when \(v\ge r\) (\(W_v=1\)). A
uniform class averages them over its band, with a cut at \(v=r\). Inside \([p,r]\), \(g_v\) is linear in \(v^2\), so a
band of half-width \(\varepsilon\) around \(\ell\) replaces \(\ell^2\) by \(\E[v^2]=\ell^2+\varepsilon^2/3\). Above \(r\),
\(t_v\) does not depend on \(v\). For the wide high class \(U[1,19]\) at \(r_1=9/2\),
\(W_H=\frac1{18}\bigl[\int_1^{9/2}\frac{v}{9/2}\,dv+(19-\frac92)\bigr]=\frac1{18}\bigl[\frac{77}{36}+\frac{29}2\bigr]
=\frac{599}{648}<1\), while \(W_H(r_0)=1\). Ownership still rises, because the rise of \(e_H\) from \(\rho\) outweighs
the fall of \(W_H\). The paper evaluates these integrals in exact rational arithmetic, with outward intervals for every
quantity that depends on \(\eu\).

\begin{intuition}[Why dispersion helps gross profit]
For each \(R\), \(v\mapsto(v-\max\{p,R\})_+\) is convex. A mean-preserving spread within a class therefore weakly
raises its gross profit, and \(B_r(\mu)\) at every belief. So the wide classes recruit, and point masses at the same
means do not. The argument ranks gross profit, not entry in every equilibrium.
\end{intuition}

\subsection{Many classes: separation and a price collision}\label{sec:10_robustness-many}
\paper{Prop.~A.6; Appendix~A.6}

With \(J\) classes the price is one number, but the posterior is a vector; the price reveals only the posterior mean
of class proceeds, \(\nu\cdot t\). Each class reaches each order through a liquidity order, so with \(\beta\le1/3\)
every ratio \(f_j(q)/f_l(q)\) lies in \([1/(1+u),1+u]\) (\Cref{sec:04_slot_beliefs-zero}). Bayes' rule,
\(\mu_j(q)/\pi_j=f_j(q)/\sum_l\pi_lf_l(q)\), then puts every order posterior in the convex set
\[
\mathcal M_\pi=\Bigl\{\nu\ge0:\ \sum_j\nu_j=1,\ \ \frac{\nu_j}{\pi_j}\le(1+u)\frac{\nu_l}{\pi_l}\ \text{for all }j,l\Bigr\},
\]
and a pooled price belief too. As in \Cref{sec:05_investor-residual}, against a fixed price and entry schedule,
\(U_j(q)=q\,e(q)\,[t_j-\mu(q)\cdot t]-k|q|\) with \(e(q)\in[\rho,1]\); this needs no price sufficiency.

\begin{resultbox}[Many classes under separation \paper{Prop.~A.6}]
Split the classes into a low group \(\mathcal L\) and a high group \(\mathcal H\), with high-group prior mass \(\bar\pi\)
and prior-weighted group means \(\bar g_{\mathcal L},\bar g_{\mathcal H}\) of \(g_{j,1}\). Suppose
\(\max_jt_{j,0}-\min_jt_{j,0}<k\), \(\sum_j\pi_jg_{j,0}<c_H\),
\[
\min_{j\in\mathcal H}t_{j,1}-\max_{\nu\in\mathcal M_\pi}\nu\cdot t_1>\frac k\rho,\qquad
\min_{\nu\in\mathcal M_\pi}\nu\cdot t_1-\max_{j\in\mathcal L}t_{j,1}>\frac k\rho,
\]
\(\bar g_{\mathcal H}>\bar g_{\mathcal L}\) and \(\bar g_{\mathcal L}+M(\bar\pi)(\bar g_{\mathcal H}-\bar g_{\mathcal L})>c_H\).
At \(F_0\), every class abstains, and entry is \(\rho\). At \(F_1\), the high group buys, the low group sells, only a
purchase recruits, and entry is \(\rho+(1-\rho)(b+\bar\pi\lambda)\).
\end{resultbox}

\emph{Proof.} At \(F_0\), every residual is at most \(\max_lt_{l,0}-\min_lt_{l,0}<k\) in absolute value, so all classes
abstain, posteriors stay at the prior, and \(\sum_j\pi_jg_{j,0}<c_H\) gives entry \(\rho\). At \(F_1\), under every
candidate, a high-group purchase earns at least \(\rho[\min_{\mathcal H}t_{l,1}-\max_\nu\nu\cdot t_1]-k>0\), and its
sale at most \(-k\); the low group is symmetric. With these orders, the class proportions within each group stay at
the prior proportions: with \(z\) the posterior mass of \(\mathcal H\), \(\mu_j=z\pi_j/\bar\pi\) on \(\mathcal H\) and
\(\mu_j=(1-z)\pi_j/(1-\bar\pi)\) on \(\mathcal L\), and \(z=(m(\bar\pi),\bar\pi,M(\bar\pi))\) by the two-class formulas.
Gross profit is \(\bar g_{\mathcal L}+z(\bar g_{\mathcal H}-\bar g_{\mathcal L})\). At \(z=\bar\pi\) it equals
\(\sum_j\pi_jg_{j,1}\le\sum_j\pi_jg_{j,0}<c_H\) (Proposition~1(a)); at \(M(\bar\pi)\) it exceeds \(c_H\). Separation
gives \(\bar t_{\mathcal H}>\bar t_{\mathcal L}\), so the three prices differ, and each reveals \(z\), hence the vector.

\emph{Four classes.} Take the benchmark slot, costs, reserve, \(r_0\) and \(r_1\), equal priors \(\reg{four_prior}\),
and point values \(\reg{four_v1}\), \(\reg{four_v2}\) (low group) and \(\reg{four_v3}\), \(\reg{four_v4}\) (high group).
Every vector in \(\mathcal M_\pi\) is \(\nu_j=\pi_jw_j/\sum_l\pi_lw_l\) with \(1\le w_j\le1+u\). The mean \(\nu\cdot t\)
is linear-fractional, hence monotone, in each weight, so its extremes lie at the 16 corners of the weight box. They
give the separation distance \(\gamma=\reg{p2_groups_gamma}\) and the strong trading margin
\(\rho\gamma-k=\reg{p2_groups_trade}\).

\paragraph{Why separation matters: a price collision.}
Take priors \((1/4,1/2,1/4)\), \(R\sim U[0,8]\), \(p=1\), \(\rho=1/4\), \(c_H=10\), \(k=1/8\), the benchmark slot, and
\(G_1=\delta_1\), \(G_2=(1-w)\delta_1+w\delta_{10}\), \(G_3=\delta_{10}\) with \(w=(1+2b_0+\lambda)/2=\reg{p2_many_weight}\).
Since \(\beta>1/4\), \(1-w=(1-\lambda)(4\beta-1)/2>0\). Class 1 sells, class 2 abstains and class 3 buys. Every class
gross profit is at most \(95/16<c_H\), so entry is \(\rho\) at every price. The auction terms give
\(t=(1,\,1+49w/16,\,65/16)\); normalize to \((0,w,1)\). After zero the posterior weights are proportional to
\((b_0/4,(b_0+\lambda)/2,b_0/4)\), and after a purchase to \((b/4,b/2,(b+\lambda)/4)\), so the normalized means are
\[
z_0=\frac{b_0(w/2+1/4)+\lambda w/2}{b_0+\lambda/2},\qquad
z_1=\frac{b(w/2+1/4)+\lambda/4}{b+\lambda/4}.
\]
Cross-multiply, cancel the common term \(bb_0(w/2+1/4)\), divide by \(\lambda\), and use \(2b+b_0+\lambda=1\):
\[
z_0=z_1\iff w\,\frac{2b+b_0+\lambda}8=\frac{3b_0}{16}+\frac b8+\frac\lambda8
\iff w=\frac{2b+3b_0+2\lambda}{2(2b+b_0+\lambda)}=\frac{1+2b_0+\lambda}2 .
\]
Both orders have the price \(\reg{p2_many_q2_price}\), but the posterior vectors differ:
\((\reg{p2_many_q1_mu0},\reg{p2_many_q1_mu1},\reg{p2_many_q1_mu2})\) after zero and
\((\reg{p2_many_q2_mu0},\reg{p2_many_q2_mu1},\reg{p2_many_q2_mu2})\) after a purchase. The chosen orders beat the best
alternatives by \(\reg{p2_many_choice0}\), \(\reg{p2_many_choice1}\) and \(\reg{p2_many_choice2}\), so this is an
equilibrium with strict choices. The collision leaves entry at \(\rho\); it refutes only the claim that the price
always reveals the whole vector. It is a knife-edge: the paper checks that a change of \(1/1000\) in \(w\), in either direction, separates the two prices and keeps every choice strict.

\subsection{Complementary private signals}\label{sec:10_robustness-signals}
\paper{Sec.~6.3; Prop.~A.7}

\begin{intuition}[Accuracy does not measure redundancy]
The investor can know the target's customers; the challenger can know its own integration plans. Two signals about
different facts each add news, even when the challenger's signal is the more accurate one.
\end{intuition}

\emph{Set-up.} Prior \(1/2\), values \(h\) and \(\ell\). The investor sees \(T\in\{+,-\}\) with accuracy \(a_T\); the
challenger sees \(Y\in\{+,-\}\) with accuracy \(d\); \(1/2<a_T,d<1\). The signals are independent given the class and
independent of the slot, the cost and \(R\). The investor orders on \(T\). Write \(\nu(x)=\Pr(T=+\mid X=x)\).

\emph{Step 1: the public belief.} The order depends on the class only through \(T\), so
\(\mu(x)=a_T\nu(x)+(1-a_T)[1-\nu(x)]=(1-a_T)+(2a_T-1)\nu(x)\). The slot bounds of \Cref{sec:04_slot_beliefs-bounds},
applied to \(T\), give \(m\le\nu\le M\). So \(\mu_-\le\mu(x)\le\mu_+\), where \(\mu_-\) and \(\mu_+\) are the values of \(\mu\) at \(\nu=m\) and
\(\nu=M\).

\emph{Step 2: the challenger adds its signal.} Bayes' rule with prior \(\mu\) gives
\(\phi_+(\mu)=d\mu/[d\mu+(1-d)(1-\mu)]\) and \(\phi_-(\mu)=(1-d)\mu/[(1-d)\mu+d(1-\mu)]<\phi_+(\mu)\). Let \(I_y\) be
high-cost preparation after \(Y=y\); since \(B\) rises, \(I_+\ge I_-\). A class-\(H\) challenger draws \(Y=+\) with
probability \(d\), so
\[
e_H=\rho+(1-\rho)[dI_++(1-d)I_-],\qquad e_L=\rho+(1-\rho)[(1-d)I_++dI_-],\qquad
0\le e_H-e_L\le(1-\rho)(2d-1).
\]

\emph{Step 3: the price still reveals \(\mu\).} Competitive prices give
\(P=t_0+\mu e_Hw_H+(1-\mu)e_Lw_L=t_0+e_Lw_L+\mu\Delta_P\), with \(\Delta_P=e_Hw_H-e_Lw_L=e_L\DT+(e_H-e_L)w_H\). Since
\(\rho\le e_L\le1\),
\[
\rho\DT\le\Delta_P\le\DT+(1-\rho)(2d-1)w_H=:D^+(r).
\]
Entry depends on the price, the signal and the cost, so two orders with one price share \(e_H\) and \(e_L\), and
\(\Delta_P\ge\rho\DT>0\) forces equal public posteriors.

\emph{Step 4: residuals.} After \(T=+\), the investor's class belief is \(a_T\), so the purchase residual is
\((a_T-\mu)\Delta_P=(2a_T-1)[1-\nu(1)]\Delta_P\). After \(T=-\), the sale residual is
\([\mu-(1-a_T)]\Delta_P=(2a_T-1)\nu(-1)\Delta_P\). With \(m\le\nu\le M\) and \(1-M=m\), every correct gross return lies in
\([(2a_T-1)\rho m\DT,\ (2a_T-1)D^+(r)]\) under every profile. A wrong-signed order earns at most \(-k\).

\begin{resultbox}[Complementary signals \paper{Prop.~A.7}]
Suppose
\[
B_{r_0}(d)<c_H<B_{r_1}\bigl(\phi_+(\mu_+)\bigr),\qquad (2a_T-1)D^+(r_0)<k<(2a_T-1)\rho m\DT(r_1).
\]
At \(r_0\), orders are zero after both signals, and entry is \(\rho\). At \(r_1\), the investor buys after \(T=+\) and
sells after \(T=-\), and entry exceeds \(\rho\). If also \(B_{r_1}(\phi_-(\mu_+))<c_H\), only \((X,Y)=(1,+)\) recruits,
and
\[
\bar e_H=\rho+(1-\rho)d(b+a_T\lambda),\qquad \bar e_L=\rho+(1-\rho)(1-d)[b+(1-a_T)\lambda],\qquad
\Ent_1=\tfrac12(\bar e_H+\bar e_L).
\]
The conditions allow \(d>a_T\).
\end{resultbox}

The trading bounds fix the orders. At \(r_0\), zero orders give \(\mu\equiv1/2\), so the challenger's belief is \(d\) or
\(1-d\), and \(B_{r_0}(d)<c_H\) keeps the high cost out. At \(r_1\), the purchase posterior is \(\mu_+\), and
\(\phi_+(\mu_+)\) recruits. In the one-pair case, a sale or zero has public posterior at most \(1/2\), and
\(B_{r_1}(d)\le B_{r_0}(d)<c_H\); a purchase has probability \(b+a_T\lambda\) in class \(H\) and \(b+(1-a_T)\lambda\) in
class \(L\).

\paragraph{The vector \(S(B)\) and an exact witness.}
Keep the benchmark slot, values, reserve and prior, and set \((\rho,c_H,k,r_0,r_1)=(17/20,357/50,3/200,11/10,23/10)\)
and \((a_T,d)=(7/10,3/4)\). Entry rises from \(\reg{p4_s_signal_r0_entry}\) to \(\reg{p4_s_signal_r1_entry}\). Without
\(Y\), the top public posterior cannot justify the high cost at either strength, and entry stays at \(\rho\): both
sources are needed. For a witness in exact fractions, set \((\lambda,\beta)=(1/3,3/10)\), so \(b=1/5\), \(u=5/3\),
\(m=3/11\) and \(M=8/11\). The five margins (private exclusion, joint recruitment, weak orders, strong orders, the
other pair) are
\[
\frac{21}{400},\qquad\frac{1363}{36800},\qquad\frac{2}{1375},\qquad\frac{603}{31625},\qquad\frac{76753}{18400}.
\]
All are positive, with \(d>a_T\).

\paragraph{The benchmark with the accuracies of \(S(B)\).}
Restore the benchmark \(\rho\), costs and strengths, and keep \((a_T,d)=(7/10,3/4)\). Three of the four conditions
fail, and each equilibrium has a mixed order. Write \(s=\lambda\sigma_+(1)\) and \(t=\lambda\sigma_-(-1)\); then
\(1-\nu(1)=b/(2b+s)\) and \(\nu(-1)=b/(2b+t)\).
\begin{itemize}[leftmargin=2em]
\item \emph{At \(r_0\).} Under every profile, \(B_{r_0}(\phi_+(\mu_-))>c_H>B_{r_0}(\phi_-(\mu_+))\): \(Y=+\) recruits at
every price and \(Y=-\) never does. So \((e_H,e_L)=(13/16,7/16)\) and
\(\Delta_P=\tfrac7{16}\cdot\tfrac1{60}+\tfrac6{16}\cdot\tfrac{33}{80}=311/1920\). The purchase gross return
\(\tfrac25\cdot\tfrac{b}{2b+s}\cdot\tfrac{311}{1920}\) falls strictly in \(s\), exceeds \(k\) at \(s=0\), and falls
short of \(k\) at \(s=\lambda\). Its unique zero-profit rate solves
\[
\frac{311b}{4800(2b+s)}=\frac1{50}\iff 2b+s=\frac{311b}{96}\iff s=\frac{119b}{96}.
\]
The sale gives the same \(t\). Each signal type trades with probability \(119/(96u)=\reg{p4_b_signal_r0_buy_rate}\),
and entry is \(\reg{p4_b_signal_r0_entry}\).
\item \emph{At \(r_1\).} \(Y=-\) never recruits. A purchase has public posterior at least \(1/2\), and
\(B_{r_1}(d)>c_H\), so \(Y=+\) recruits after it, with \(\Delta_P=\tfrac7{16}\cdot\tfrac23+\tfrac6{16}\cdot\tfrac98=137/192\).
The purchase earns at least \(\tfrac25m\tfrac{137}{192}-k>0\), so \(s=\lambda\). After a sale, \(Y=+\) recruits
exactly when \(\phi_+(\mu(-1))\ge\tau(r_1)=141/200\), that is \(\mu(-1)\ge47/106\), \(\nu(-1)\ge19/53\), \(t\le15b/19\).
On that range the sale earns at least \(\tfrac25\cdot\tfrac{19}{53}\cdot\rho\DT(r_1)=19/795>k\), so the investor
would sell with probability one; since \(\lambda>15b/19\), no equilibrium has \(t\le15b/19\). Above that range, entry
after a sale is \(\rho\), and the net sale payoff \(\tfrac25\cdot\tfrac{b}{2b+t}\cdot\tfrac14\cdot\tfrac23-\tfrac1{50}
=\tfrac{b}{15(2b+t)}-\tfrac1{50}\) falls strictly in \(t\). Its unique zero is \(t=4b/3\in(15b/19,\lambda)\): the sale
probability is \(4/(3u)=\reg{p4_b_signal_r1_sell_rate}\). After zero, \(Y=+\) also recruits, so only a sale blocks
recruitment after \(Y=+\). A sale has probability \(b+(1-a_T)t\) in class \(H\) and \(b+a_Tt\) in class \(L\), so
\[
\Ent_1=\rho+\frac{1-\rho}2\bigl\{d[1-b-(1-a_T)t]+(1-d)(1-b-a_Tt)\bigr\}=\frac58-\frac{23b}{40}=\reg{p4_b_signal_r1_entry}.
\]
\end{itemize}
Entry falls in this comparison. Rule IM lets the investor mix at zero profit, and the strict monotonicity of each
return makes these tuples unique. Under rule I0, which forces abstention at zero profit, neither strength has an
equilibrium. Without \(Y\), entry is \(\rho\) at both strengths, with zero orders at \(r_0\) and rates \(s=t=4b/3\) at
\(r_1\), again under rule IM.

\paragraph{The benchmark itself.}
At the benchmark the conditions hold with \((a_T,d)=(\reg{sig_comp_aT},\reg{sig_comp_d})\). Entry rises to
\(\reg{p4_comp_b_r1_entry}\), and it stays at \(\rho\) without \(Y\). The conditions cannot hold there with \(d\ge a_T\).
Write \(z=2a_T-1\). The strong bound needs \(z>z_*=k/(m\rho\DT(r_1))\). If \(d\ge a_T\), then \(2d-1\ge z\), and the
weak gross bound is \(zD^+(r_0)\ge z_*[\DT(r_0)+(1-\rho)z_*w_H(r_0)]>k\), which violates the weak bound. This limits
these sufficient conditions at the benchmark only: \(S(B)\) has \(d>a_T\).

\subsection{Order sizes and the Laplace comparison}\label{sec:10_robustness-sizes}
\paper{Sec.~2.3, fn.~6}

A liquidity trader places each of the orders \(\{-1,0,1\}\) with positive probability, so every order is on path, and
Bayes' rule defines every price. A menu of continuous sizes needs a convention for the sizes that no class uses,
unless liquidity flow covers every size. The Laplace version took the second route. With additive Laplace noise of
scale \(s\) (density proportional to \(\eu^{-|z|/s}\)) and full orders \(\pm1\), the likelihood ratio of the flow \(x\)
is
\[
\frac{f(x-1)}{f(x+1)}=\exp\Bigl(\frac{|x+1|-|x-1|}{s}\Bigr)\le\eu^{2/s},
\]
since \(|x+1|-|x-1|\le2\). So the two set-ups have the same posterior bounds when \(1+u=\eu^{2/s}\), that is
\(s=2/\log(1+u)\). The benchmark \(1+u=\eu\) corresponds to \(s=2\). The match concerns \(m\) and \(M\), not
the experiment or the equilibrium. In the slot, class \(H\) reaches the top posterior with probability \(b+\lambda\)
(a purchase); in the Laplace model, it reaches the top plateau \(x\ge1\) with probability \(1/2\). These masses enter
entry even when the bounds agree.

\begin{panelnote}[What changed from the Laplace version]
The slot removes the derivative argument over order sizes: you compare three payoffs directly. The old
logistic-noise extension has no counterpart. Cost atoms, the moderate value gap and complementary signals still
support the reversal, and many classes add a new point: one price need not reveal a whole posterior vector.
\end{panelnote}
